\subtitle{\bf Aula 6 - Análise Descendente}
\begin{frame}[t,plain]
\titlepage
\end{frame}

\begin{frame}{Objetivos}
\protect\hypertarget{objetivos-1}{}
\begin{itemize}
\item
  Apresentar a técnica de análise sintática descendente recursiva.
\item
  Mostrar como combinadores de parsing implementam essa estratégia em
  Haskell.
\end{itemize}
\end{frame}



\begin{frame}{Análise Descendente Recursiva}
\protect\hypertarget{anuxe1lise-descendente-recursiva}{}
\begin{itemize}

\item
  Técnica para construir analisadores sintáticos como um
  \textbf{conjunto de funções}
\item
  Uma função por não-terminal da gramática
\end{itemize}
\end{frame}

\begin{frame}{Análise Descendente Recursiva}
\protect\hypertarget{anuxe1lise-descendente-recursiva-1}{}
\begin{itemize}

\item
  A chamada de funções reflete a estrutura das produções
\item
  Em Haskell: abordagem elegante com \textbf{combinadores de parsing}
\item
  Um parser é tratado como um \textbf{valor de primeira classe}
\end{itemize}
\end{frame}

\begin{frame}[fragile]{O Tipo \texttt{Parser}}
\protect\hypertarget{o-tipo-parser}{}
\begin{lstlisting}[language=Haskell]
newtype Parser s a =
  Parser { runParser :: [s] -> [(a, [s])] }
\end{lstlisting}

\begin{itemize}

\item
  \texttt{s}: tipo dos símbolos de entrada (ex.:
  \texttt{Char}, \texttt{Token})
\item
  \texttt{a}: tipo do resultado (ex.:
  \texttt{Exp}, \texttt{[Token]})
\end{itemize}
\end{frame}

\begin{frame}[fragile]{O Tipo \texttt{Parser}}
\protect\hypertarget{o-tipo-parser-1}{}
\begin{itemize}

\item
  Retorna \textbf{lista de pares} (resultado, entrada restante)
\item
  Lista vazia \texttt{[]} significa \textbf{falha}
\item
  Lista com vários elementos significa \textbf{ambiguidade}
\end{itemize}
\end{frame}

\begin{frame}[fragile]{Resultados}
\protect\hypertarget{resultados}{}
\begin{itemize}

\item
  Sucesso sem não determinismo.

  \begin{itemize}

  \item
    Apenas um resultado.
  \end{itemize}
\end{itemize}

\begin{lstlisting}[language=Haskell]
runParser (symbol 'a') "abc"
-- -> [('a', "bc")]    -- sucesso: consome 'a', resta "bc"
\end{lstlisting}
\end{frame}

\begin{frame}[fragile]{Resultados}
\protect\hypertarget{resultados-1}{}
\begin{itemize}

\item
  Falha: Lista vazia de resultados
\end{itemize}

\begin{lstlisting}[language=Haskell]
runParser (symbol 'x') "abc"
-- -> []               -- falha: 'a' ≠ 'x'
\end{lstlisting}
\end{frame}

\begin{frame}[fragile]{Resultados}
\protect\hypertarget{resultados-2}{}
\begin{itemize}

\item
  Mais de um resultado: não determinismo

  \begin{itemize}

  \item
    Pode resultar em backtracking.
  \end{itemize}
\end{itemize}

\begin{lstlisting}[language=Haskell]
runParser (item <|> item) "abc"
-- -> [('a', "bc"), ('a', "bc")]
\end{lstlisting}
\end{frame}

\begin{frame}{Resultados}
\protect\hypertarget{resultados-3}{}
\begin{itemize}

\item
  Parsers determinísticos retornam no máximo um par
\item
  A lista permite parsers \textbf{não-determinísticos} de forma elegante
\end{itemize}
\end{frame}

\hypertarget{instuxe2ncias-de-classes-de-tipos}{%
\section{Instâncias de Classes de
Tipos}\label{instuxe2ncias-de-classes-de-tipos}}

\begin{frame}[fragile]{Functor}
\protect\hypertarget{functor}{}
\begin{itemize}

\item
  Um parser é um \texttt{Functor}

  \begin{itemize}

  \item
    Responsável por executar ações.
  \item
    Ex. Construir AST.
  \end{itemize}
\end{itemize}

\begin{lstlisting}[language=Haskell]
instance Functor (Parser s) where
  fmap f p = Parser (\s ->
    [(f x, s') | (x, s') <- runParser p s])
\end{lstlisting}
\end{frame}

\begin{frame}[fragile]{Functor}
\protect\hypertarget{functor-1}{}
\begin{itemize}

\item
  Transforma o resultado sem modificar o que é consumido
\item
  Exemplo: \texttt{EInt <$> naturalParser} aplica
  \texttt{EInt} ao resultado
\end{itemize}
\end{frame}

\begin{frame}[fragile]{Applicative}
\protect\hypertarget{applicative}{}
\begin{itemize}

\item
  Concatenação (sequência) de parsers.
\end{itemize}

\begin{lstlisting}[language=Haskell]
instance Applicative (Parser s) where
  pure a    = Parser (\ts -> [(a, ts)])
  p1 <*> p2 = Parser (\s ->
    [(f x, s2) | (f, s1) <- runParser p1 s
               , (x, s2) <- runParser p2 s1])
\end{lstlisting}
\end{frame}

\begin{frame}[fragile]{Applicative}
\protect\hypertarget{applicative-1}{}
\begin{itemize}

\item
  \texttt{pure a}: sempre tem sucesso sem consumir
  entrada
\item
  \texttt{p1 <*> p2}: executa
  \texttt{p1} e depois \texttt{p2} com
  a entrada restante
\end{itemize}
\end{frame}

\begin{frame}[fragile]{Alternative}
\protect\hypertarget{alternative}{}
\begin{itemize}

\item
  Escolha entre dois parsers.
\end{itemize}

\begin{lstlisting}[language=Haskell]
instance Alternative (Parser t) where
  empty     = Parser (\_ -> [])
  p1 <|> p2 = Parser (\s ->
    runParser p1 s ++ runParser p2 s)
\end{lstlisting}
\end{frame}

\begin{frame}[fragile]{Alternative}
\protect\hypertarget{alternative-1}{}
\begin{itemize}

\item
  \texttt{empty}: sempre falha
\item
  \texttt{p1 <|> p2}: tenta ambas as alternativas
  (concatena resultados)
\item
  Devido à avaliação lazy, \texttt{p2} só é avaliado se
  \texttt{p1} falhar
\end{itemize}
\end{frame}

\begin{frame}[fragile]{Monad}
\protect\hypertarget{monad}{}
\begin{itemize}

\item
  Permite a definição de parsers que dependem do resultado de um parser
  anterior.
\end{itemize}

\begin{lstlisting}[language=Haskell]
instance Monad (Parser t) where
  return = pure
  p >>= f = Parser (\ts -> concat
    [ runParser (f a) cs'
    | (a, cs') <- runParser p ts ])
\end{lstlisting}
\end{frame}

\begin{frame}[fragile]{Monad}
\protect\hypertarget{monad-1}{}
\begin{itemize}

\item
  Permite que o segundo parser \textbf{dependa do resultado} do primeiro
\item
  Habilita notação \texttt{do} para parsers mais
  legíveis
\end{itemize}
\end{frame}

\hypertarget{combinadores-primitivos}{%
\section{Combinadores Primitivos}\label{combinadores-primitivos}}

\begin{frame}[fragile]{Combinador \texttt{item}}
\protect\hypertarget{combinador-item}{}
\begin{lstlisting}[language=Haskell]
-- Consome um símbolo qualquer
item :: Parser t t
item = Parser (\ts -> case ts of
  []     -> []
  (c:cs) -> [(c, cs)])
\end{lstlisting}
\end{frame}

\begin{frame}[fragile]{Combinador \texttt{sat}}
\protect\hypertarget{combinador-sat}{}
\begin{lstlisting}[language=Haskell]
-- Consome um símbolo se satisfaz o predicado
sat :: (t -> Bool) -> Parser t t
sat p = do
  t <- item
  if p t then return t else mzero
\end{lstlisting}
\end{frame}

\begin{frame}[fragile]{Combinador \texttt{symbol}}
\protect\hypertarget{combinador-symbol}{}
\begin{lstlisting}[language=Haskell]
symbol :: Eq s => s -> Parser s s
symbol c = sat (c ==)
\end{lstlisting}
\end{frame}

\begin{frame}[fragile]{Combinador \texttt{token}}
\protect\hypertarget{combinador-token}{}
\begin{lstlisting}[language=Haskell]
token :: Eq s => [s] -> Parser s [s]
token = mapM symbol
\end{lstlisting}
\end{frame}

\begin{frame}[fragile]{Combinador \texttt{digit}}
\protect\hypertarget{combinador-digit}{}
\begin{lstlisting}[language=Haskell]
digit :: Parser Char Int
digit = f <$> sat isDigit
  where f c = ord c - ord '0'
\end{lstlisting}
\end{frame}

\begin{frame}[fragile]{Combinador \texttt{natural}}
\protect\hypertarget{combinador-natural}{}
\begin{lstlisting}[language=Haskell]
natural :: Parser Char Int
natural = foldl (\a b -> a*10 + b) 0 <$> many1 digit
\end{lstlisting}
\end{frame}

\hypertarget{combinadores-de-alto-nuxedvel}{%
\section{Combinadores de Alto
Nível}\label{combinadores-de-alto-nuxedvel}}

\begin{frame}[fragile]{Repetição}
\protect\hypertarget{repetiuxe7uxe3o}{}
\begin{lstlisting}[language=Haskell]
-- Um ou mais
many :: Parser s a -> Parser s [a]
many p = (:) <$> p <*> many p <|> pure []
\end{lstlisting}
\end{frame}

\begin{frame}[fragile]{Repetição}
\protect\hypertarget{repetiuxe7uxe3o-1}{}
\begin{lstlisting}[language=Haskell]
-- Um ou mais
many1 :: Parser s a -> Parser s [a]
many1 p = (:) <$> p <*> many p
\end{lstlisting}
\end{frame}

\begin{frame}[fragile]{Opção}
\protect\hypertarget{opuxe7uxe3o}{}
\begin{lstlisting}
-- Zero ou um
option :: Parser s a -> a -> Parser s a
option p v = p <|> succeed v
\end{lstlisting}
\end{frame}

\begin{frame}[fragile]{Delimitadores}
\protect\hypertarget{delimitadores}{}
\begin{lstlisting}[language=Haskell]
-- Reconhece p entre dois delimitadores
pack :: Parser s a -> Parser s b -> Parser s c -> Parser s b
pack p r q = pi32 <$> p <*> r <*> q
\end{lstlisting}
\end{frame}

\begin{frame}[fragile]{Delimitadores}
\protect\hypertarget{delimitadores-1}{}
\begin{lstlisting}[language=Haskell]
-- Expressão entre parênteses:
parens p = pack (symbol '(') p (symbol ')')
\end{lstlisting}
\end{frame}

\begin{frame}[fragile]{Separadores}
\protect\hypertarget{separadores}{}
\begin{lstlisting}[language=Haskell]
-- Lista separada por delimitador
listOf :: Parser s a -> Parser s b -> Parser s [a]
listOf p s = list <$> p <*> many (pi22 <$> s <*> p)
\end{lstlisting}
\end{frame}

\begin{frame}[fragile]{Repetição Gulosa}
\protect\hypertarget{repetiuxe7uxe3o-gulosa}{}
\begin{lstlisting}[language=Haskell]
determ :: Parser s b -> Parser s b
determ p = Parser (\ts -> case runParser p ts of
  []    -> []
  (x:_) -> [x])

greedy  :: Parser s b -> Parser s [b]
greedy  = determ . many
\end{lstlisting}
\end{frame}

\begin{frame}[fragile]{Operadores}
\protect\hypertarget{operadores}{}
\begin{itemize}
\item
  Para gramáticas da forma \(E \to E \oplus T \mid T\):
\item
  \texttt{chainl}: associatividade \textbf{à esquerda}
  --- \texttt{1+2+3} ->
  \texttt{(1+2)+3}
\end{itemize}

\begin{lstlisting}[language=Haskell]
chainl :: Parser s a -> Parser s (a -> a -> a) -> Parser s a
chainl pe po = h <$> pe <*> many (j <$> po <*> pe)
  where j op x = (`op` x)
        h x fs = foldl (flip ($)) x fs
\end{lstlisting}
\end{frame}

\begin{frame}[fragile]{Operadores}
\protect\hypertarget{operadores-1}{}
\begin{itemize}
\item
  Para gramáticas da forma \(E \to T \oplus E \mid T\):
\item
  \texttt{chainr}: associatividade \textbf{à direita}
  --- \texttt{1+2+3} ->
  \texttt{1+(2+3)}
\end{itemize}

\begin{lstlisting}[language=Haskell]
chainr :: Parser s a -> Parser s (a -> a -> a) -> Parser s a
chainr pe po = h <$> many (j <$> pe <*> po) <*> pe
  where j x op = (x `op`)
        h fs x = foldr ($) x fs
\end{lstlisting}
\end{frame}

\hypertarget{parser-de-expressuxf5es}{%
\section{Parser de Expressões}\label{parser-de-expressuxf5es}}

\begin{frame}[fragile]{Gramática e Tokens}
\protect\hypertarget{gramuxe1tica-e-tokens}{}
\begin{lstlisting}[language=Haskell]
data Token = Id String | Number Int | Add | Mult
           | LParen | RParen deriving (Eq, Show)

data Expr = Var String | Lit Int
          | Expr :+: Expr | Expr :*: Expr
\end{lstlisting}
\end{frame}

\begin{frame}[fragile]{O Combinador \texttt{gen}}
\protect\hypertarget{o-combinador-gen}{}
\begin{lstlisting}[language=Haskell]
type Op s a = (s, a -> a -> a)

gen :: Eq s => [Op s a] -> Parser s a -> Parser s a
gen ops p = chainl p (choice (map f ops))
  where f (s, c) = const c <$> symbol s
\end{lstlisting}

\begin{itemize}

\item
  \texttt{gen ops p} reconhece expressões com
  operadores de \texttt{ops} e operandos de
  \texttt{p}
\item
  Associatividade à esquerda via \texttt{chainl}
\end{itemize}
\end{frame}

\hypertarget{parser-de-expressuxf5es-1}{%
\section{Parser de Expressões}\label{parser-de-expressuxf5es-1}}

\begin{frame}[fragile]{Parser de Expressões}
\protect\hypertarget{parser-de-expressuxf5es-2}{}
\begin{itemize}

\item
  \(E \to E + T \,|\, T\)
\end{itemize}

\begin{lstlisting}[language=Haskell]
exprParser :: Parser Token Expr
exprParser = addtable `gen` termParser
  where addtable = [(Add, (:+:))]
\end{lstlisting}
\end{frame}

\begin{frame}[fragile]{Parser de Expressões}
\protect\hypertarget{parser-de-expressuxf5es-3}{}
\begin{itemize}

\item
  \(T \to T * F \,|\, F\)
\end{itemize}

\begin{lstlisting}[language=Haskell]
termParser :: Parser Token Expr
termParser = multable `gen` factParser
  where multable = [(Mult, (:*:))]
\end{lstlisting}
\end{frame}

\begin{frame}[fragile]{Parser de Expressões}
\protect\hypertarget{parser-de-expressuxf5es-4}{}
\begin{lstlisting}[language=Haskell]
factParser :: Parser Token Expr
factParser = numParser <|> varParser
          <|> pack (symbol LParen)
                   exprParser
                   (symbol RParen)
\end{lstlisting}
\end{frame}

\hypertarget{megaparsec}{%
\section{Megaparsec}\label{megaparsec}}

\begin{frame}{Por que Megaparsec?}
\protect\hypertarget{por-que-megaparsec}{}
\begin{itemize}

\item
  A biblioteca didática é para ensino; \textbf{Megaparsec} é para
  produção
\item
  Mensagens de erro detalhadas: linha, coluna, contexto, o que era
  esperado
\end{itemize}
\end{frame}

\begin{frame}[fragile]{Por que Megaparsec?}
\protect\hypertarget{por-que-megaparsec-1}{}
\begin{itemize}

\item
  Melhor desempenho com \emph{committed choice}
\item
  Combinador \texttt{makeExprParser} para hierarquias
  de precedência
\end{itemize}
\end{frame}

\begin{frame}[fragile]{O Tipo Parser no Megaparsec}
\protect\hypertarget{o-tipo-parser-no-megaparsec}{}
\begin{lstlisting}[language=Haskell]
type Parser a = Parsec Void String a
-- Parsec tipo-de-erro tipo-de-entrada tipo-de-resultado
\end{lstlisting}

\begin{itemize}

\item
  \texttt{Void}: sem tipo de erro customizado (usa os
  padrões da biblioteca)
\item
  \texttt{String}: entrada como lista de caracteres
\end{itemize}
\end{frame}

\begin{frame}[fragile]{Tratamento de Espaços}
\protect\hypertarget{tratamento-de-espauxe7os}{}
\begin{itemize}

\item
  Remover espaços e lidar com comentários
\end{itemize}

\begin{lstlisting}[language=Haskell]
sc :: Parser ()
sc = L.space space1 lineCmnt blockCmnt
  where lineCmnt  = L.skipLineComment "//"
        blockCmnt = L.skipBlockComment "/*" "*/"
\end{lstlisting}
\end{frame}

\begin{frame}[fragile]{Tokens}
\protect\hypertarget{tokens}{}
\begin{lstlisting}
lexeme :: Parser a -> Parser a
lexeme = L.lexeme sc  -- consome espaço após o token

symbol :: String -> Parser String
symbol = L.symbol sc

parens :: Parser a -> Parser a
parens = between (symbol "(") (symbol ")")
\end{lstlisting}
\end{frame}

\begin{frame}[fragile]{Operadores com
\texttt{makeExprParser}}
\protect\hypertarget{operadores-com-makeexprparser}{}
\begin{lstlisting}[language=Haskell]
opTable :: [[Operator Parser Exp]]
opTable
  = [ [ infixL (:*:) "*" ]   -- maior precedência
    , [ infixL (:+:) "+" ]   -- menor precedência
    ]
  where infixL op sym = InfixL $ op <$ symbol sym
\end{lstlisting}
\end{frame}

\begin{frame}[fragile]{Operadores}
\protect\hypertarget{operadores-2}{}
\begin{itemize}

\item
  Lista externa: do \textbf{maior} para o \textbf{menor} nível de
  precedência
\item
  Lista interna: operadores de mesma precedência
\item
  \texttt{InfixL}: associatividade à esquerda;
  \texttt{InfixR}: à direita
\end{itemize}
\end{frame}

\begin{frame}[fragile]{Parser com
\texttt{makeExprParser}}
\protect\hypertarget{parser-com-makeexprparser}{}
\begin{lstlisting}[language=Haskell]
termP :: Parser Exp
termP = parens expP <|> EInt <$> int

expP :: Parser Exp
expP = makeExprParser termP opTable
\end{lstlisting}
\end{frame}

\begin{frame}[fragile]{Parser com
\texttt{makeExprParser}}
\protect\hypertarget{parser-com-makeexprparser-1}{}
\begin{itemize}

\item
  \texttt{makeExprParser} constrói automaticamente toda
  a hierarquia de precedências
\item
  Sem necessidade de codificar manualmente cada nível como função
  separada
\end{itemize}
\end{frame}

\hypertarget{conclusuxe3o}{%
\section{Conclusão}\label{conclusuxe3o}}

\begin{frame}[fragile]{Conclusão}
\protect\hypertarget{conclusuxe3o-1}{}
\begin{itemize}

\item
  \textbf{Combinadores de parsing} tratam parsers como valores de
  primeira classe
\item
  O tipo \texttt{Parser s a} retorna lista de pares
  (resultado, entrada restante)
\item
  Instâncias de \texttt{Functor},
  \texttt{Applicative},
  \texttt{Alternative}, \texttt{Monad}
  dão expressividade
\end{itemize}
\end{frame}

\begin{frame}[fragile]{Conclusão}
\protect\hypertarget{conclusuxe3o-2}{}
\begin{itemize}

\item
  Primitivos: \texttt{item},
  \texttt{sat}, \texttt{symbol};
  combinadores: \texttt{many},
  \texttt{option}, \texttt{pack},
  \texttt{chainl}
\item
  O combinador \texttt{gen} +
  \texttt{chainl} resolve associatividade sem
  transformar a gramática
\end{itemize}
\end{frame}

\begin{frame}{Conclusão}
\protect\hypertarget{conclusuxe3o-3}{}
\begin{itemize}

\item
  \textbf{Megaparsec}: mesmos princípios, porém com uma implementação
  eficiente.
\end{itemize}
\end{frame}
